\begin{abstract}

Este trabalho investiga a aplicação de modelos de regressão para a predição do desempenho de um sistema de dessalinização por destilação por membranas do tipo V-AGMD (\textit{Vacuum-Enhanced Air Gap Membrane Distillation}), com o objetivo de estimar o fluxo de permeado e as temperaturas de saída das correntes a partir de variáveis operacionais do processo. A relevância do estudo está associada à complexidade dos fenômenos de transferência de calor e massa envolvidos, que dificultam a obtenção de modelos precisos por abordagens exclusivamente físicas ou empíricas. Embora o fluxo de permeado seja adotado como variável prioritária no processo de seleção dos modelos, as temperaturas de saída também são consideradas por fornecerem informações complementares sobre o comportamento térmico do sistema e por estarem associadas à avaliação de indicadores energéticos, como o consumo específico de energia (SEC) e o ganho de saída (GOR).

O problema é formulado como uma tarefa de regressão supervisionada multissaída, utilizando dados experimentais obtidos em diferentes regimes de operação. A validação dos modelos é realizada por meio de validação cruzada com separação por grupos experimentais. São comparadas diferentes famílias de modelos, incluindo regressões lineares regularizadas, modelos baseados em árvores e redes neurais artificiais, sendo a seleção conduzida com base no erro de predição do fluxo de permeado e em critérios de simplicidade do modelo. Os resultados são analisados em comparação com um modelo físico simplificado, permitindo discutir as vantagens e limitações das abordagens baseadas em dados, bem como evidenciar o potencial de estratégias híbridas para a melhoria da capacidade de generalização na modelagem do processo.

\end{abstract}